% !TeX root = ../MAgent_Theory_2.0.tex
\chapter{Sufficient States for Smooth KL Evolution}
\label{ch:smooth-closure}

This chapter works with the reciprocal law potential of \cref{def:dyn-potential} and its unit-mobility Fisher flow; its conclusions concern that potential, not the per-agent or global free energies of \cref{ch:dynamics}. The question is which information a coarse description of an agent must keep when an exact reduction is wanted while agents remain different. We first identify the conditional information that enters a retained force, then construct an exact state completion made of general probability laws for one architecture: fixed independent private reference factors, and receiver reconstruction channels that discard the sender's private coordinate. These channels are additional constitutive data within the potential, not a consequence of an arbitrary law-dependent transport.

In this architecture the private sector is passive. Each private conditional relaxes along an explicit exponential arc toward its reference, independently of every public law, row, and other agent, and it acts on the public flow only as open-loop forcing through its conditional KL profile. The completion is therefore a solvable baseline. Its compression is the equimeasurability quotient of each private log density under its declared reference: it forgets which private point carries a log-density value, so it identifies two conditionals exactly when their log densities are equimeasurable under the reference. On a finite alphabet this requires two distinct sets of atoms of equal mass, which is weaker than a symmetry of the reference. The completion is not minimal for the public flow (\cref{smooth:cgf-completion}), and the per-agent completion fails once a channel transmits the sender's private data (\cref{smooth:channel-variants}). The construction selects no group of agents, so clauses (F2), (F5), and (F6) of \cref{inh:recursive-formation} are not addressed; its private sector evolves identically under both nulls and under the frozen-environment control of \cref{inh:null-flow}, because the conditional equations contain no peer term.

\section{The force that survives marginalization}

Let $Q$ be a positive law on a standard Borel space, let $C$ be a measurable statistic, and write $\mu=C_\#Q$. Suppose a differentiable functional has variational derivative $h_Q$ and Fisher velocity
\begin{equation}
 \dot Q=-Q(h_Q-\E_Qh_Q).
 \label{smooth:fine-score}
\end{equation}
All identities below are immediate for finite positive tables. On general spaces they need a dominated differentiable path, square-integrable scores, and integrability that permits the displayed conditional expectations and differentiations. Statements about flows on general spaces are made under the hypotheses of \cref{smooth:hypotheses}, which give unique solutions in bounded log-density charts (\cref{smooth:wellposed}); they do not assert existence on every space of probability measures.

\propositionheading{Conditional mean force}{smooth:mean-force-criterion}
The retained velocity is
\begin{equation}
 \dot\mu=-\mu\left(\E_Q[h_Q\mid C]-\E_Qh_Q\right).
 \label{smooth:mean-force}
\end{equation}
Consequently, an autonomous velocity for $\mu$ exists on an admitted class precisely when the centered conditional mean on the right is determined by $\mu$ on that class. If, in addition, the admitted class is invariant under the fine flow and the resulting retained equation $\dot\mu=V(\mu)$ has unique solutions, then two fine trajectories in the class with the same initial marginal have the same marginal at all common times.
\begin{proof}
For every bounded measurable $f$, integrate $f(C)$ against \cref{smooth:fine-score} and condition on $C$. This gives the weak form of \cref{smooth:mean-force}. Two laws with the same marginal have the same marginal velocity if and only if their centered conditional forces agree almost everywhere under that marginal. Invariance keeps both fine trajectories in the class, so each marginal solves $\dot\mu=V(\mu)$, and uniqueness identifies them.
\end{proof}

Equation \cref{smooth:mean-force} is the class-mean aggregation of a replicator equation over types, familiar from the Price equation of selection \citep{Price1970}. For $Q(dz,dv)=\mu(dz)\kappa_z(dv)$ and a fixed reference $P(dz,dv)=\nu(dz)\pi_z(dv)$, the local KL force illustrates the missing information:
\begin{equation}
 \E_Q[\log(dQ/dP)\mid z]
 =\log\frac{d\mu}{d\nu}(z)+H(z),\qquad
 H(z)=\KL(\kappa_z\Vert\pi_z).
 \label{smooth:hidden-profile}
\end{equation}
Thus a scalar total conditional loss does not usually determine the retained force. Its dependence on $z$ matters. The Fisher metric makes the same separation: if a score is $f(z,v)=f_1(z)+f_2(z,v)$, with $\E_\mu f_1=0$ and $\E_{\kappa_z}f_2=0$, then
\begin{equation}
 g_Q(fQ,fQ)=\int f_1(z)^2\mu(dz)
       +\int\left[\int f_2(z,v)^2\kappa_z(dv)\right]\mu(dz).
 \label{smooth:fisher-split}
\end{equation}
The cross term vanishes by conditional centering. This calculation is consistent with the statistical behavior of Fisher information under sufficient statistics \citep{AyJostLeSchwachhofer2015}. Dynamical sufficiency asks a different question: it preserves the specified evolution and responses, and need not preserve inference about every discarded parameter.

\corollaryheading{Pushforward as a Fisher submersion}{smooth:submersion}
On a regular positive class, the differential of $Q\mapsto C_\#Q$ sends the tangent $fQ$ to $\E_Q[f\mid C]\mu$. Its kernel consists of the scores with zero conditional mean given $C$, and it is isometric on scores that are functions of $C$. If the functional has the form $\mathcal E=\mathcal E_c\circ C_\#$, the retained marginal follows the Fisher gradient flow of $\mathcal E_c$ at the same time.
\begin{proof}
For bounded $g$, $\int g(C)f dQ=\int g\E_Q[f\mid C]d\mu$, which identifies the pushforward tangent. With $f_1=\E_Q[f\mid C]$ and $f_2=f-f_1$, \cref{smooth:fisher-split} gives the kernel and the isometry. If $\mathcal E=\mathcal E_c\circ C_\#$, the chain rule gives $h_Q=h_c\circ C$ up to a constant, so \cref{smooth:mean-force} becomes $\dot\mu=-\mu(h_c-\E_\mu h_c)$.
\end{proof}
This is the equality case of the monotonicity of Fisher information under statistics \citep{AyJostLeSchwachhofer2018}; the content of the present chapter is the case in which $h_Q$ is not a function of $C$.

\section{Conditional reconstruction within the original potential}

For each agent take $\mathsf K_i=\mathsf Z_i\times\mathsf V_i$ and $\mathsf M_i=\mathsf W_i\times\mathsf U_i$, all standard Borel. The coordinates $z,w$ are public in the sense that the channels below read them; the coordinates $v,u$ are private in the sense that no channel reads the sender's value of them, and the declared probes read them only through the marks of \cref{sec:smooth-records}. Use public densities against fixed reference probabilities $\rho_i$ on $\mathsf Z_i$ and $\varrho_i$ on $\mathsf W_i$, so conditional versions share common full-measure public sets. Disintegrate
\begin{equation}
 q_i(dz,dv)=\mu_i(dz)\kappa_{i,z}(dv),\qquad
 s_i(dw,du)=\sigma_i(dw)\varkappa_{i,w}(du).
 \label{smooth:disintegrations}
\end{equation}
Neither conditional need be independent of its retained coordinate. Choose fixed private reference probabilities $A_i^q,A_i^s$ and a fixed normalized public reference
\begin{equation}
 R_i^c(dz,dw)=r_i^c(dw)G_i^c(w,dz),\qquad
 R_i=R_i^c\otimes A_i^q\otimes A_i^s,
 \label{smooth:reference}
\end{equation}
with the product reordered into state and model coordinates. The evaluator $G_i^c$ can couple $\mu_i$ and $\sigma_i$. Independence of the private reference factors is an assumption of this construction, not a restriction of the general theory.

Let $B_{ij}^q[\boldsymbol\mu]:\mathsf Z_j\rightsquigarrow\mathsf Z_i$ be normalized channels, differentiable in the laws in the sense required by \cref{smooth:hypotheses} on general spaces. Define the full channels by
\begin{equation}
 \mathcal K_{ij}^q[\mathbf q]((z_j,v_j),d(z_i,v_i))
 =B_{ij}^q[\boldsymbol\mu](z_j,dz_i)\kappa_{i,z_i}(dv_i).
 \label{smooth:shield}
\end{equation}
The model channel uses $B_{ij}^s[\boldsymbol\sigma]$ and $\varkappa_{i,w_i}$ in the same way. Source sets may exclude self edges; if self comparisons are included, use the primitive identity convention separately. Public channels may be informative and nonconstant functions of the retained laws. For example, on the public space $\{-1,1\}$, the channel $B_\mu(z,\{z'\})=(1+\vartheta(\mu)zz')/2$, with $\vartheta(\mu)=(1+(\E_\mu Z)^2)/4$, is strictly positive, transmits source information, and varies with the law. Public/private maps and reference factors are declared interface data; carrying them through a bimeasurable change of presentation preserves these constructions.

\propositionheading{Exact cancellation of private peer information}{smooth:peer-cancellation-rule}
Where the divergences are defined,
\begin{equation}
 \begin{aligned}
 q_j\mathcal K_{ij}^q&=(\mu_jB_{ij}^q)\kappa_i,\\
 \KL(q_i\Vert q_j\mathcal K_{ij}^q)
 &=\KL(\mu_i\Vert\mu_jB_{ij}^q).
 \end{aligned}
 \label{smooth:peer-cancellation}
\end{equation}
This is an identity on the full admitted class, including correlated $q_i$. Its derivative includes the sender, receiver, and channel derivatives.
\begin{proof}
Integration over the source conditional gives the first identity. The compared laws have the same receiver conditional $\kappa_{i,z}$, so their density ratio depends only on $z$. Integrating its logarithm proves the second identity. Differentiating the resulting functional identity on any regular path proves the derivative statement. Holding $\kappa_i$ fixed in the denominator leaves the force at the frozen state unchanged, because the derivative of the divergence in the receiver conditional of the compared law vanishes there; it violates the identity away from that state, changes the second variation, and, if the snapshot is held while the laws move, changes the force.
\end{proof}

Let $\mathcal E_c$ denote exactly \cref{eq:dyn-potential} applied to $\mu_i,\sigma_i,R_i^c,B_{ij}^q,B_{ij}^s$, with the same rows, source priors, and temperatures. Define
\begin{equation}
 \begin{aligned}
 \ell_{i,z}^q&=\log(d\kappa_{i,z}/dA_i^q),&
 H_i^q(z)&=\int\ell_{i,z}^q d\kappa_{i,z},\\
 \ell_{i,w}^s&=\log(d\varkappa_{i,w}/dA_i^s),&
 H_i^s(w)&=\int\ell_{i,w}^s d\varkappa_{i,w}.
 \end{aligned}
 \label{smooth:conditional-entropies}
\end{equation}
On finite public alphabets with strictly positive $r_i^c$ and $G_i^c$, twice continuously differentiable channel entries, and positive compared laws on the interior, (H2) and (H3) below hold, and it suffices that each private log density be bounded. On general public spaces we use the following hypotheses.

\hypothesisheading{Regularity on general spaces}{smooth:hypotheses}
Use the bounded log-density chart for every public law: $\mu_i=e^{u_i}\rho_i/\E_{\rho_i}e^{u_i}$ with $u_i$ in $L^\infty_0(\rho_i)=\{u\in L^\infty(\rho_i):\E_{\rho_i}u=0\}$, the same for $\sigma_i$ against $\varrho_i$, and log-odds coordinates for interior rows.
\begin{enumerate}
\item[(H1)] The initial private log densities are bounded uniformly in the public coordinate: at time zero,
\begin{equation}
 \operatorname*{ess\ sup}_{(z,v)}|\ell^q_{i,z}(v)|\leq M<\infty
 \label{smooth:h1-bound}
\end{equation}
with respect to $\rho_i\otimes A_i^q$, and likewise in the model sector.
\item[(H2)] The public reference has bounded log densities: $G_i^c(w,dz)=g_i^c(w,z)\rho_i(dz)$ and $r_i^c=\tilde r_i\varrho_i$, with $\log g_i^c\in L^\infty(\varrho_i\otimes\rho_i)$ and $\log\tilde r_i\in L^\infty(\varrho_i)$.
\item[(H3)] On bounded subsets of the product of the charts, the public first variations $h_i^q=\delta\mathcal E_c/\delta\mu_i$ and $h_i^s=\delta\mathcal E_c/\delta\sigma_i$, including every interaction role and channel derivative, are Lipschitz into $L^\infty(\rho_i)$ and $L^\infty(\varrho_i)$, and the costs $D_{ij}^q,D_{ij}^s$ are Lipschitz. Identity channels between agents sharing a public reference space satisfy this, as do channels with kernel densities $b_{ij}[\boldsymbol\mu]$ whose logarithms are bounded on bounded subsets, provided $\boldsymbol\mu\mapsto b_{ij}[\boldsymbol\mu]$ is continuously differentiable into $L^\infty(\rho_j\otimes\rho_i)$ with a derivative represented by bounded kernels, $Db_{ij}[\boldsymbol\mu](\delta\mu_k)=\int\beta_{ijk}[\boldsymbol\mu](z',\cdot,\cdot)\delta\mu_k(dz')$ with $\beta_{ijk}\in L^\infty(\rho_k\otimes\rho_j\otimes\rho_i)$ locally Lipschitz in $\boldsymbol\mu$. Channels that depend on the laws through moments $\int g d\mu_k$ qualify when $g$ is bounded.
\end{enumerate}
Positivity of the public densities alone does not suffice: a normalized, strictly positive, law-dependent channel whose first variation contains an unbounded function of $z$ gives a public force outside $L^\infty$, and the chart vector field is then undefined. Continuous differentiability into $L^\infty$ alone does not exclude this: on $(0,1]$ with Lebesgue reference, the kernel density $1+\varsigma(y)\cos(2\pi z')\cos(2\pi z)$, with $y=\int\log(1/z)\mu(dz)$ and a smooth $|\varsigma|\leq1/2$, is a smooth function of the chart coordinate into $L^\infty$ because $\log(1/z)$ is integrable, but the first variation of a peer cost through it contains $\varsigma'(y)\log(1/z)$ times a factor that is nonzero in general. A bound on each $\ell_{i,z}$ separately does not suffice either: a family of bounded conditionals whose bounds grow as $z$ approaches a boundary point can make the public log density unbounded for every positive time. The class (H1) excludes conditionals whose density vanishes on a set of positive reference measure or is unbounded, such as Gaussian conditionals against a Gaussian reference; \cref{smooth:exponential-family} treats those separately.

\theoremheading{Exact marginal and conditional dynamics}{smooth:split-dynamics}
For the architecture above the original energy is
\begin{equation}
 \mathcal E=\mathcal E_c+
 \sum_i\int H_i^q(z)\mu_i(dz)+
 \sum_i\int H_i^s(w)\sigma_i(dw).
 \label{smooth:energy-split}
\end{equation}
Write $h_i^q=\delta\mathcal E_c/\delta\mu_i$ and $h_i^s=\delta\mathcal E_c/\delta\sigma_i$, including every interaction role. The full unit-mobility Fisher dynamics gives
\begin{equation}
 \begin{aligned}
 \dot\mu_i&=-\mu_i\left[h_i^q+H_i^q-
                  \E_{\mu_i}(h_i^q+H_i^q)\right],\\
 \partial_t\ell_{i,z}^q&=H_i^q(z)-\ell_{i,z}^q.
 \end{aligned}
 \label{smooth:split-flow}
\end{equation}
The equations for $\sigma_i,\ell_i^s$ are obtained by exchanging sectors. Source rows evolve by their original Fisher equations with the public costs, or may be exactly optimized as in \cref{prop:dyn-gibbs}.
\begin{proof}
Apply the relative-entropy chain rule to the local reference in \cref{smooth:reference} and use \cref{smooth:peer-cancellation}. This proves \cref{smooth:energy-split}. Up to a constant independent of $(z,v)$, the full derivative with respect to $q_i$ is $h_i^q(z)+\ell_{i,z}^q(v)$. This can be checked directly from the log-density ratio in the local KL; every peer contribution is a functional of $\mu$ alone. Conditional averaging gives the first equation of \cref{smooth:split-flow}. Subtract its logarithmic derivative from the full logarithmic derivative of $q_i=\mu_i\kappa_i$. The public score and global centering cancel, leaving the second equation. The private terms contain no source rows, so their variations do not change the row equations.
\end{proof}

The two row choices share stationary points but are different flows. At frozen costs the Fisher row equation relaxes log-ratio contrasts at the rate $T_q$ (\cref{smooth:row-ratio}), so exactly optimized rows are the limit of infinite row mobility, and the two choices differ most in their transients when $T_q$ is small.

The private equation in \cref{smooth:split-flow} contains no public law, no row, and no other agent. Each private conditional is solved in closed form, independently of the public dynamics, by \cref{smooth:conditional-solution} below: the exponential arc from $\kappa_{i,z}(0)$ to $A_i^q$. The public flow sees the private sector only through the function $H_i^q(z,t)$, which is fixed by the initial conditionals. The completion constructed next is therefore the public flow of $\mathcal E_c$ driven by the precomputable time-dependent potential $\int H_i^q(\cdot,t)d\mu_i$. Private information never influences an interaction in this architecture. The dependence of $H_i^q$ on $z$ is the only private contribution to the retained belief velocity, and it cannot generally be discarded.

\section{An exact law-valued completion}

\definitionheading{Conditional log-density spectrum}{smooth:spectrum-definition}
For a private reference $A$ and conditional density $\kappa_z$, define
\begin{equation}
 \xi_z=(\ell_z)_\#A,\qquad
 \ell_z=\log(d\kappa_z/dA).
 \label{smooth:spectrum}
\end{equation}
This probability law records how much reference mass carries each log-density value. It forgets which private point carries that value. It obeys
\begin{equation}
 \int e^r\xi_z(dr)=1,\qquad
 H(z)=\int r e^r\xi_z(dr).
 \label{smooth:spectrum-normalization}
\end{equation}

The spectrum is the law under $A$ of the log-likelihood ratio of the dichotomy $(\kappa_z,A)$. For mutually absolutely continuous pairs, equal spectra are equivalent to Blackwell equivalence of the two dichotomies \citep{Blackwell1953,Torgersen1991}. Two conditionals have the same spectrum exactly when their log densities are equimeasurable under $A$. The compression is therefore a property of the declared reference, not of the dynamics. On a finite private alphabet $\xi_z(\{r\})=A(\{\ell_z=r\})$, so if no two distinct sets of atoms of $A$ have equal mass, equal spectra force equal level sets and hence equal conditionals: the retention is injective. With $A=(0.6,0.4)$, for example, the conditionals $(3/4,1/4)$ and $(1/4,3/4)$ have the different spectra with atoms $\log(5/4),\log(5/8)$ and $\log(5/12),\log(15/8)$ at masses $0.6,0.4$. Conversely, if two distinct sets of atoms have equal mass, their set differences $S,S'$ are disjoint, nonempty, and of equal mass, and exchanging two distinct density values placed on $S$ and $S'$ gives distinct conditionals with equal spectra. Noninjectivity therefore requires a reference with equal-mass sets of atoms, such as a uniform finite reference, or a reference with a nonzero atomless part, in which disjoint sets of equal positive mass always exist. The log-Laplace transform of the spectrum is
\begin{equation}
 \Lambda_z(b)=\log\int e^{br}\xi_z(dr)=(b-1)D_b(\kappa_z\Vert A),
 \label{smooth:renyi}
\end{equation}
where $D_b$ is the R\'enyi divergence of order $b$ \citep{vanErvenHarremoes2014}. For every $\kappa_z$ mutually absolutely continuous with $A$ and $0<b<1$, H\"older's inequality gives $\int e^{br}\xi_z(dr)=\int(d\kappa_z/dA)^bdA\in(0,1]$, and this function is analytic in the strip $0<\operatorname{Re}b<1$. A law is determined by its moment generating function on any nondegenerate interval on which it is finite, so $\Lambda_z$ on any subinterval of $(0,1)$ determines $\xi_z$; (H1) is not needed for this.

\theoremheading{Closed conditional-spectrum evolution}{smooth:spectrum-flow}
For \cref{smooth:split-dynamics}, the velocity of the state consisting of all $\mu_i,\sigma_i$, conditional spectra $\xi_i^q,\xi_i^s$, and source rows is a function of that state. Each spectrum satisfies, in weak form,
\begin{equation}
 \frac{d}{dt}\int f(r)\xi_z(dr)
 =\int f'(r)(H(z)-r)\xi_z(dr),
 \label{smooth:spectrum-weak}
\end{equation}
for continuously differentiable test functions on the initial range $I=[m_0,M_0]$ of the log density, which contains the range at every later time. For a time interval $t$ starting at zero, its exact solution is
\begin{equation}
 \begin{aligned}
 a&=e^{-t},& Z_z(a)&=\int e^{ar}\xi_z^0(dr),\\
 \xi_z^t&=(r\mapsto ar-\log Z_z(a))_\#\xi_z^0,
 \end{aligned}
 \label{smooth:spectrum-solution}
\end{equation}
and it is the only narrowly continuous curve of probability laws supported in $I$ that solves \cref{smooth:spectrum-weak} with $H(z)$ given by \cref{smooth:spectrum-normalization} from the same initial law with $\int e^r\xi^0_z(dr)=1$. The entropy profile obeys
\begin{equation}
 \dot H(z)=-\operatorname{Var}_{e^r\xi_z}(r).
 \label{smooth:entropy-rate}
\end{equation}
\begin{proof}
The reference $A$ is fixed. Differentiate $\int f(\ell_z(v))A(dv)$ and substitute the conditional equation in \cref{smooth:split-flow}; this proves \cref{smooth:spectrum-weak}. Direct differentiation shows that
\begin{equation}
 \frac{d\kappa_z(t)}{dA}(v)
 =\frac{\exp(a\ell_z^0(v))}{Z_z(a)}
 \label{smooth:conditional-solution}
\end{equation}
solves the normalized conditional equation; this is the explicit solution of the Fisher--Rao gradient flow of a relative entropy, the exponential arc toward the reference \citep{ChenEtAl2023}. Pushing $A$ through its logarithm gives \cref{smooth:spectrum-solution}. For the range, H\"older's inequality gives $Z_z(a)\leq1$, and $e^{ar}\geq e^re^{-(1-a)M_0}$ for $r\leq M_0$ gives $Z_z(a)\geq e^{-(1-a)M_0}$; since $m_0\leq0\leq M_0$ by normalization, $a\ell_z^0-\log Z_z(a)$ takes values in $[m_0,M_0]$, and the width of its range is $a(M_0-m_0)$.

For uniqueness, let $\xi_t$ be a curve in the stated class and put $H(t)=\int re^r\xi_t(dr)$, which is continuous. Then $\xi_t$ solves the linear continuity equation with the velocity field $H(t)-r$, which is Lipschitz in $r$, and continuity equations with such fields have unique narrowly continuous solutions \citep[Sec.~8.1]{AmbrosioGigliSavare2008}. Hence $\xi_t=(X_t)_\#\xi_0$ with $X_t(r)=e^{-t}r+c(t)$ and $c(t)=\int_0^te^{-(t-s)}H(s)ds$. Taking $f(r)=e^r$ in \cref{smooth:spectrum-weak} gives $\dot N=H(N-1)$ for $N=\int e^r\xi_t(dr)$, so $N\equiv1$, which forces $e^{c(t)}Z(a)=1$ and gives \cref{smooth:spectrum-solution}.

Setting $f(r)=re^r$ in the weak equation gives
$\dot H=H(H+1)-\int(r^2+r)e^r\xi_z(dr)$, which is \cref{smooth:entropy-rate}. The probability $e^r\xi_z(dr)$ is normalized by \cref{smooth:spectrum-normalization}. All private contributions to \cref{smooth:split-flow} are now functions of the retained spectra. The remaining scores and row equations depend only on retained public data, so the retained velocity is a function of the retained state.
\end{proof}

That the retained state determines its own future additionally requires uniqueness of the public flow forced by $H_i(z,t)$. For finite public alphabets with twice continuously differentiable channels and the positivity hypotheses of \cref{thm:dyn-finite-flow}, its Picard--Lindel\"of argument applies to this forced system. On general public spaces the following proposition supplies it.

\propositionheading{Well-posedness and exact lift on general spaces}{smooth:wellposed}
Assume \cref{smooth:hypotheses}. Then for every initial state in the charts:
\begin{enumerate}
\item[(i)] each private log density keeps its values in its initial essential range, $0\leq H_i^q(z,t)\leq H_i^q(z,0)\leq M$, and $|\partial_tH_i^q(z,t)|\leq M^2$ uniformly in $z$, and likewise in the model sector;
\item[(ii)] the retained system of public laws, rows, and spectra has a unique maximal solution on an interval $[0,T^*)$, with public laws and rows continuously differentiable into the charts and spectra given by \cref{smooth:spectrum-solution};
\item[(iii)] the fine Fisher system has a unique maximal continuously differentiable solution in the charts $L^\infty_0(\rho_i\otimes A_i^q)$ and $L^\infty_0(\varrho_i\otimes A_i^s)$ on the same interval, and it equals the lift $q_i(t)=\mu_i(t)\kappa_i(t)$ with $\kappa_i(t)$ given by \cref{smooth:conditional-solution}, and likewise in the model sector;
\item[(iv)] if $R$ sends a fine state to its public laws, spectra, and rows, then $R\circ\Phi_t=\Phi_t^c\circ R$ on $[0,T^*)$, where $\Phi^c$ is the solution map of (ii).
\end{enumerate}
\begin{proof}
(i) The range statement is in the proof of \cref{smooth:spectrum-flow}, and it holds uniformly in $z$ by (H1). The function $H$ is a relative entropy, hence nonnegative, and it is nonincreasing by \cref{smooth:entropy-rate}. A law on an interval of width at most $2M$ has variance at most $M^2$.

(ii) In the chart the public equation reads $\dot u_i=-(\phi_i-\E_{\rho_i}\phi_i)$ with $\phi_i=h_i^q+H_i^q(\cdot,t)$, because then $\partial_t\log(d\mu_i/d\rho_i)=-(\phi_i-\E_{\mu_i}\phi_i)$. This is the chart flow of \cref{def:dyn-chart-flow} for $\mathcal E_c$, with the time-dependent term $H_i^q(\cdot,t)$ added to the first variation. By (i), $t\mapsto H_i^q(\cdot,t)$ is Lipschitz into $L^\infty(\rho_i)$, and (H2) and (H3) make $h_i^q$ Lipschitz on bounded sets; the row block is a finite-dimensional equation whose costs are Lipschitz on bounded sets, and Gibbs rows are smooth functions of the costs. The proof of \cref{prop:dyn-chart-flow}, applied to the system augmented by the time coordinate, gives a unique maximal solution, with the alternative that it leaves every bounded set as $t\to T^*$ if $T^*<\infty$.

(iii) In the fine chart, the first variation of the energy with respect to $q_i$ is $\varphi_i(z,v)=h_i^q(z)+\ell_{i,z}(v)$ up to a constant, by the proof of \cref{smooth:split-dynamics}. The public marginal map $l\mapsto\log\int e^{l(\cdot,v)}A_i^q(dv)$ is smooth from $L^\infty(\rho_i\otimes A_i^q)$ to $L^\infty(\rho_i)$, so the fine first variation is Lipschitz on bounded sets, and \cref{prop:dyn-chart-flow} gives uniqueness of fine solutions. For the lift, $\partial_t\log q_i=-(h_i^q+H_i^q-\E_{\mu_i}(h_i^q+H_i^q))+(H_i^q-\ell_i)=-(\varphi_i-\E_{q_i}\varphi_i)$, because $\E_{q_i}\varphi_i=\E_{\mu_i}h_i^q+\E_{\mu_i}H_i^q$. The lift stays bounded on compact subintervals of $[0,T^*)$ by (i), so it is the fine solution there. Conversely, the projection of a fine solution solves the retained system, so a fine solution bounded beyond $T^*$ would contradict the alternative in (ii).

(iv) Two fine states with the same retained value have the same public laws, rows, and spectra, hence the same forcing $H(t)$ and, by (ii), the same retained solution; by (iii) their projections are retained solutions.
\end{proof}

Under these hypotheses the retained state is therefore autonomous on the common maximal interval. For law-dependent public channels, global existence, $T^*=\infty$, requires further estimates on the public sector, which are not given here. For fixed public channels it holds.

\corollaryheading{Global existence for fixed public channels}{smooth:global}
Assume \cref{smooth:hypotheses} and, in addition, that every public channel $B_{ij}^q,B_{ij}^s$ is fixed and either an identity between agents sharing a public reference or a kernel whose density lies in $[c,C_K]$ with $0<c\leq C_K<\infty$, and that the rows are fixed or follow their Fisher equations. Then $T^*=\infty$ in \cref{smooth:wellposed}.
\begin{proof}
The public system is the chart flow of $\mathcal E_c$ with the first variation of each public law increased by $H_i^q(\cdot,t)$ or $H_i^s(\cdot,t)$, which lies in $[0,M]$ by \cref{smooth:wellposed}(i). Run the proof of \cref{prop:dyn-chart-global} for this system. In the lower bound the first variation increases by at most $M$ and its mean does not decrease, so the constant $A$ of \cref{eq:dyn-chart-global-constants} is replaced by $A+M$. In the upper bound $\mathcal E_c(t)$ is at most the full energy \cref{smooth:energy-split} at time $t$, which is nonincreasing along the fine flow, and the forcing adds at most $M$ to the mean of the first variation, so the bound holds with $B$ computed from the full initial energy and increased by $M$. The public log densities therefore stay bounded on every finite interval, and the continuation alternative of \cref{smooth:wellposed}(ii) gives $T^*=\infty$; by (iii) the fine solution is global as well.
\end{proof}

This is a configuration retention. Its spectra can be infinite-dimensional, and there is no claim that they give a fixed finite number of parameters. If represented as a sample statistic, the map $v\mapsto\ell_z(v)$ depends on the current law; its motion is exactly the term differentiated in \cref{smooth:spectrum-weak}. It must not be frozen while the law changes. The inherited spectral evolution is not a Fisher--Rao gradient flow on probability laws over $(z,r)$. Such a flow $\dot Q=-Q(h-\E_Qh)$ keeps $Q_t$ mutually absolutely continuous with $Q_0$, whereas the retained conditional $e^r\xi^t_z(dr)$ is supported on $a\operatorname{supp}\xi_z^0-\log Z_z(a)$, which moves unless $\kappa_z=A$. The spectral evolution is a continuity equation in $r$ with velocity $H(z)-r$.

\corollaryheading{Different internal conditionals with the original public flow}{smooth:equal-spectra}
Suppose that, for each agent and each sector separately, the initial spectrum is independent of the public coordinate. Then this property persists, and the public laws and source rows solve the original coarse KL/Fisher system generated by $\mathcal E_c$ on the common maximal interval.
\begin{proof}
The same initial spectrum gives the same solution of \cref{smooth:spectrum-solution} at every public coordinate. Consequently $H_i^q(z)$ is constant in $z$, and $H_i^s(w)$ is constant in $w$. Each constant disappears under its own marginal centering in \cref{smooth:split-flow}.
\end{proof}

Equal spectra do not require equal conditionals when the reference is uniform. With a private reference uniform on two points, the conditionals $(3/4,1/4)$ and $(1/4,3/4)$ share
\begin{equation}
 \xi=\tfrac12\delta_{\log(3/2)}+
      \tfrac12\delta_{\log(1/2)}.
 \label{smooth:swapped-spectrum}
\end{equation}
A public bit may select either conditional, producing correlation between public and private variables. Another law with the same public marginal may use the first conditional at both public values and be a product. When the original public flow from the common initial public state is unique, as it is for this finite public alphabet by \cref{thm:dyn-finite-flow} and on general spaces under \cref{smooth:hypotheses}, both have the same public trajectory in this architecture, although their internal laws differ. Different agents may have different public laws and different spectra. No equality between agents is imposed.

This class is invariant, but it is not reached for the first time at finite interior time. For any $a>0$, the spectral update in \cref{smooth:spectrum-solution} is invertible on its image. Indeed, if $\xi^t$ is given, the initial log value is recovered by the map
\begin{equation}
 y\longmapsto\frac{y}{a}
 -\log\int e^{r/a}\xi^t(dr).
 \label{smooth:spectrum-inverse}
\end{equation}
Substituting \cref{smooth:spectrum-solution} and using $\int e^r\xi^0(dr)=1$ proves the assertion. Equal spectra at a finite time therefore had to be equal initially. The same inversion applied to \cref{smooth:conditional-solution} proves that distinct conditionals cannot first become equal at finite time. Neither fact forbids an autonomous quotient that keeps the unequal spectra.

\propositionheading{A coarser completion for the public flow}{smooth:cgf-completion}
In the setting of \cref{smooth:wellposed}, the public laws and rows are determined, on the common maximal interval, by their initial values together with the classes of the maps $z\mapsto\Lambda^q_{i,z}|_{(0,1]}$ and $w\mapsto\Lambda^s_{i,w}|_{(0,1]}$ of \cref{smooth:renyi} modulo functions of $b$ alone. These classes evolve autonomously by
\begin{equation}
 \Lambda^t_z(b)=\Lambda^0_z(ab)-b\Lambda^0_z(a),\qquad a=e^{-t}.
 \label{smooth:cgf-flow}
\end{equation}
Conversely, two private configurations with the same initial public laws and rows whose public trajectories agree on some interval $[0,\varepsilon)$ have the same classes, so the class is exactly the private information that the public flow needs and reveals. The class is strictly coarser than the spectrum: if $\mathsf V=\mathsf V_1\times\mathsf V_2$ with a product reference, the conditionals $\kappa^{(1)}_z\otimes\psi$ and $\kappa^{(1)}_z\otimes\psi'$ with $\psi\neq\psi'$ fixed have different spectra at every $z$ whenever $\psi$ and $\psi'$ have different spectra, but the same class.
\begin{proof}
Substituting \cref{smooth:spectrum-solution} into the definition gives \cref{smooth:cgf-flow}. If $\Lambda^0_z=\widetilde\Lambda^0_z+k(b)$ with $k$ independent of $z$, then \cref{smooth:cgf-flow} adds $k(ab)-bk(a)$, again independent of $z$, so the flow acts on classes. The forcing is $H(z,t)=\partial_b\Lambda^t_z(1)$, which the class determines up to a function of $t$ alone, and such a function disappears under marginal centering in \cref{smooth:split-flow}. Uniqueness in \cref{smooth:wellposed} gives the conclusion. For the example, the log-Laplace transform of a product conditional against a product reference is the sum of the transforms of the factors, so the two maps differ by the $z$-independent function $\Lambda_\psi-\Lambda_{\psi'}$, while the spectra are the convolutions of a common factor with the different spectra of $\psi$ and $\psi'$.

For the converse, equal public trajectories have equal public velocities, so $H(\cdot,t)-\widetilde H(\cdot,t)=k(t)$ almost everywhere for each $t$ in the interval, with one null set for all $t$ because both sides are continuous in $t$. With $\Delta_z=\Lambda^0_z-\widetilde\Lambda^0_z$, $a=e^{-t}$, and $H(z,t)=\partial_b\Lambda^t_z(1)=a^2\frac{d}{da}\bigl(\Lambda^0_z(a)/a\bigr)$, this reads $\frac{d}{da}(\Delta_z(a)/a)=k(t(a))/a^2$ on $(e^{-\varepsilon},1]$. Since $\Delta_z(1)=0$, $\Delta_z(a)=aK(a)$ there with $K$ independent of $z$. Each $\Lambda^0_z$ is real analytic on $(0,1)$, so $\Delta_z-\Delta_{z_0}$ vanishes on $(0,1]$, and the two classes coincide.
\end{proof}
Beyond the class, the full spectrum is needed for responses that read the conditional spectra themselves, such as $f$-divergence readouts $\int f(e^r)\xi_z(dr)$, and its marked version below for records that read private data. \Cref{smooth:equal-spectra} is the case in which the classes are trivial.

\remarkheading{Exponential families and finite-parameter private closure}{smooth:exponential-family}
H\"older's inequality gives $0<Z_z(a)\leq1$ for every $\kappa_z(0)$ mutually absolutely continuous with $A$, so \cref{smooth:conditional-solution} defines the private solution without (H1). If $\kappa_z(0)\propto\exp(\theta_z\cdot T)A$ for a statistic $T$ with log-partition function $\psi$, then $\kappa_z(t)\propto\exp(e^{-t}\theta_z\cdot T)A$: the natural parameters contract linearly, and $H(z,t)=a\theta_z\cdot\nabla\psi(a\theta_z)-\psi(a\theta_z)$ with $a=e^{-t}$. The private sector then closes on finitely many parameters per public value, for instance on Gaussian conditionals relative to a Gaussian reference. Such families fall outside (H1) when $T$ is unbounded, and the public part then needs its own justification; for the public equation alone, a bound on $\operatorname*{ess\ sup}_zH(z,0)$ suffices for the forcing to be bounded, because $H$ is nonincreasing in $t$. Nothing here makes the public laws finite-dimensional.

\remarkheading{Which channels permit the completion}{smooth:channel-variants}
Receiver reconstruction is not what produces closure. If every channel in a row is replaced by $B_{ij}^q(z_j,dz_i)A_i^q(dv_i)$, a fixed reconstruction independent of state and source, the chain rule gives $D_{ij}^q=\KL(\mu_i\Vert\mu_jB_{ij}^q)+\int H_i^q d\mu_i$. The added term is the same for every source, so it cancels from Gibbs rows and from the row velocities, the private energy doubles, and the conditional equation becomes $\partial_t\ell_{i,z}^q=2(H_i^q(z)-\ell_{i,z}^q)$; the spectrum still closes. What matters is that the channel discards the sender's private coordinate. A channel that transmits it fails. Take a trivial public space, $\mathsf V=\{1,2\}$ with uniform $A$, two agents attending to each other through the identity channel, and $\kappa_1=(3/4,1/4)$. With $\kappa_2=\kappa_1$ the peer costs vanish, whereas with $\kappa_2=(1/4,3/4)$, which has the same spectrum, each peer cost is $\tfrac12\log3$ and the peer forces move $\kappa_1$ toward $\kappa_2$. The spectral velocity of agent $1$ then differs between two configurations with the same per-agent spectra, so those spectra are not a closed state. For identity private channels on a common $(\mathsf V,A)$, every peer integrand is a function of the vector $(\ell_{k}(v))_k$ at one private point, which suggests the joint law of that vector under $A$ as the retained state. For identity public and private channels with base-only references this joint law is closed (\cref{prop:boundary-joint-spectrum}); its closure with law-dependent public channels, and its compression beyond equimeasurable rearrangements of the private value vector, are not established here.

\section{Boundary records and feedback}
\label{sec:smooth-records}

Public-only response kernels descend immediately: the channels $B_{ij}^q,B_{ij}^s$ and evaluator $G_i^c$ are determined by the retained public laws. A private-sensitive response needs more information. In \cref{smooth:swapped-spectrum}, let the probabilities of record $1$ at the two private points be $(3/4,1/4)$. The two conditionals yield record probabilities $5/8$ and $3/8$. After record $1$, their private posteriors are $(9/10,1/10)$ and $(1/2,1/2)$. Thus an unmarked spectrum can fail both the boundary test and the next-update test.

There is an exact completion for a specified family of records. Fix a measurable mark $T:\mathsf V\to\mathsf D$, with $\mathsf D$ standard Borel, through which every allowed private readout factors. Retain
\begin{equation}
 \eta_z=(\ell_z,T)_\#A.
 \label{smooth:marked-spectrum}
\end{equation}
The correlation between the log-density value and the mark is retained. For a normalized record kernel with strictly positive density $L(o\mid z,d)$, jointly Borel in $(o,z,d)$, assume the following integrals are finite and nonzero and each updated law stays in the admitted regular domain. A bounded log likelihood is a sufficient local condition for the conditional log-density bounds. Put
\begin{equation}
 \begin{aligned}
 N_z(o)&=\int e^r L(o\mid z,d)\eta_z(dr,dd),\\
 N(o)&=\int N_z(o)\mu(dz).
 \end{aligned}
 \label{smooth:record-normalizers}
\end{equation}

\propositionheading{Exact marked feedback}{smooth:feedback}
The record density is $N(o)$. Exact conditioning of this sector gives the closed updates
\begin{equation}
 \begin{aligned}
 \mu^+(dz)&=\frac{N_z(o)}{N(o)}\mu(dz),\\
 \eta_z^+&=\bigl((r,d)\mapsto
       (r+\log L(o\mid z,d)-\log N_z(o),d)\bigr)_\#\eta_z.
 \end{aligned}
 \label{smooth:feedback-update}
\end{equation}
Between records, \cref{smooth:spectrum-weak,smooth:spectrum-solution} transport the first coordinate of $\eta_z$ and leave its mark fixed. Hence, for the stage kernels \cref{eq:dyn-record-interrupted} of the record-interrupted dynamics with these records and updates, the retention $R_{\rm ret}$ to the public laws, marked spectra, and rows satisfies $(R_{\rm ret}\times\mathrm{id})_\#J_k(x;\cdot)=J^c_k(R_{\rm ret}(x);\cdot)$ for a retained stage kernel $J^c_k$, at every state whose flow interval lies inside the maximal interval of \cref{smooth:wellposed}, so the marked state determines every finite adaptive history law of these records, feedback updates, and original smooth-flow intervals up to the stopping time of \cref{def:adaptive-history}.
\begin{proof}
Since $d\kappa_z=e^{\ell_z}dA$, integration of the likelihood gives \cref{smooth:record-normalizers}. Bayes' rule gives the displayed marginal update and
$d\kappa_z^+/dA=e^{\ell_z}L(o\mid z,T)/N_z(o)$. Taking the logarithm and pushing $A$ gives the marked update. The fine reference remains fixed. The conditional equation after the update is therefore the same equation with this new conditional law. Inductively, the inherited flow gives the next retained state, the inherited record kernel gives the same next-record law, and the inherited update gives the same posterior retained state, which is the stated identity of stage kernels. Actions selected from retained histories by Borel policies in the sense of \cref{def:adaptive-history} preserve this induction, because each stage kernel then depends on the past only through retained quantities. The history laws are defined when the flow map $(t,x)\mapsto\Phi_t(x)$ is jointly Borel on its domain, which \cref{def:dyn-record-interrupted} requires and which is part of the hypothesis on these charts. Any reference or likelihood parameters changed by records must also be retained and updated by a rule that depends only on this state.
\end{proof}

The composite of flow intervals and these updates is the record-interrupted dynamics of \cref{def:dyn-record-interrupted}, with record kernel $N(o)$ at the pre-record state and the update \cref{smooth:feedback-update}, stopped at the maximal existence time. The potential contains no record term and is not monotone across records. Records act only at their times: no rule changing the references is declared here, and between records each conditional relaxes toward its fixed reference by \cref{smooth:conditional-solution}, so the effect of a record on the private sector decays.

The update is sectorwise. A product of separate belief and model likelihoods preserves the recognition restriction $q_i\otimes s_i$. A likelihood jointly coupling the two sectors need not preserve that product; exact joint conditioning would require a joint recognition state, or a separately declared projection. Marked feedback can also destroy equality of unmarked spectra across public coordinates. The augmented marked state still closes, but the simpler conclusion of \cref{smooth:equal-spectra} then need not hold.

For an explicit retained sample law, use
\begin{equation}
 Q^c(dz,dr,dd)=\mu(dz)e^r\eta_z(dr,dd).
 \label{smooth:marked-sample-law}
\end{equation}
It is normalized and reproduces every declared marked record. It is the pushforward under the current, law-dependent statistic $(z,v)\mapsto(z,\ell_z(v),T(v))$. Different assignments of density values to private points within a mark can give the same retained state. If the permitted records distinguish every private point, the mark may have to retain that point, eliminating this compression. The boundary protocol therefore fixes what can be forgotten. Predictive-state constructions likewise make sufficiency relative to future observations \citep{ShaliziCrutchfield2001}; the formulas here specify the state and update for this particular KL architecture.

\section{What is still required for formation}

Positive source temperature does not give exact isolation by itself. For optimized rows, every allowed finite-cost edge has positive weight. For a dynamic row, set
\begin{equation}
 \theta_{jk}=\log\frac{\beta_j/\pi_j}{\beta_k/\pi_k}.
\end{equation}
The Fisher row equation from \cref{eq:dyn-potential} gives
\begin{equation}
 \dot\theta_{jk}=-(D_j-D_k)-T_q\theta_{jk}.
 \label{smooth:row-ratio}
\end{equation}
Multiplying by $e^{T_qt}$ and integrating shows that finite cost differences keep every such ratio finite on a regular finite interval. Initially positive weights therefore stay positive there. Uniformly bounded cost differences also bound these ratios for all time and prevent asymptotic disappearance of an allowed weight. A formed collective may nevertheless keep nonzero boundary interactions; isolation is not required for inheritance.

The construction above compresses private conditional information of each agent separately, and only up to the equimeasurability quotient of its declared reference. It does not combine several distinct public agent laws into one smaller collective state. Retaining their complete labeled tuple and calling it a parent would leave that aggregation problem untouched; whether a retention of such a tuple compresses is decided by clause (F4) of \cref{inh:recursive-formation}, which requires a noninjective retention whose fibers some admitted child-level probe distinguishes. Nor does recognizing a persistent source-preference group establish a compatible collective retention. The distinction between a coarse observable and an autonomous macroscopic evolution is also expressed by flow intertwining in \citet{BarnettSeth2023}.

\remarkheading{Aggregation with conditional information retained}{smooth:formation-open}
The construction of this chapter selects no group and involves no grouping of agents, and its private sector is unchanged by deleting every peer cost. A recursive formation that uses the conditional-spectrum state as part of its retention is an instance of \cref{inh:formation-open}. For identity public and private channels with base-only references, private conditionals do feed back into public interactions, and the joint conditional spectrum is an exact completion (\cref{prop:boundary-joint-spectrum}); it compresses only through equimeasurable rearrangements of the private value vector and keeps every public child state. What remains open is a coupled architecture under the reciprocal potential, with law-dependent public channels or rows that read private data, whose exact completion compresses beyond such rearrangements or forgets distinct public child states.
